% Dinâmica: 2 blocos e 1 força 
% Faro, seg 08 dez 2025 20:10:14  
% Written by Tordar 
% pstricks2pdf.sh 2blocos_1forca
\documentclass{article}
\pagestyle{empty}

\usepackage{pst-all} 
\usepackage{pstricks-add}
\usepackage{pst-slpe}
\usepackage{graphicx}
\input{apnum}

% Bloco
\newcommand{\bloco}[5][]{% #1=options, #2=cx, #3=cy, #4=totalWidth, #5=height
  % Left frame (half of total width, full height)
  \psframe[linewidth=1pt,linecolor=gray,fillstyle=slope,slopebegin=white,slopeend=gray,slopesteps=15]
    (!#2 #4 2 div sub #3 #5 2 div sub)                     % lower-left of total area
    (!#2 #3 #5 2 div add)                                  % upper-right of left frame (center x, top)  
  % Right frame (half of total width, full height)
  \psframe[linewidth=1pt,linecolor=gray,fillstyle=slope,slopebegin=gray,slopeend=white,slopesteps=15]
    (!#2 #3 #5 2 div sub)                                  % lower-left of right frame (center x, bottom)
    (!#2 #4 2 div add #3 #5 2 div add)                     % upper-right of total area
  \psframe[#1](!#2 #4 2 div sub #3 #5 2 div sub)(!#2 #4 2 div add #3 #5 2 div add) 
}

\begin{document}

% Altura e largura do bloco 1: 
\evaldef\wbi{3}
\evaldef\hbi{2}
% Altura e largura do bloco 2: 
\evaldef\wbii{3}
\evaldef\hbii{1.5}
% Coordenadas do primeiro bloco: 
\evaldef\xbi{3}
\evaldef\ybi{2+\hbi/2}
% Coordenadas do segundo bloco: 
\evaldef\xbii{\xbi+\wbi/2+\wbii/2}
\evaldef\ybii{2+\hbii/2}
% Coordenadas auxiliares:
\evaldef\xmi{\xbi}
\evaldef\ymi{\ybi+0.75*\hbi}
\evaldef\xmii{\xbii}
\evaldef\ymii{\ybii+0.75*\hbii}
% Força:
\evaldef\xfe{\xbi-\wbi/2}
\evaldef\xfs{\xfe-1.5}
\evaldef\xf{(\xfs+\xfe)/2}
\evaldef\yf{\ybi+0.5}
% Aceleração:
\evaldef\xas{\xbii+0.6*\wbii}
\evaldef\xae{\xas+1.5}
\evaldef\xa{(\xas+\xae)/2}
\evaldef\ya{\ybii+0.5}

% Let's go:
\begin{pspicture}(0,0)(10,10)
% Chão: 
\psframe[linewidth=0pt,linecolor=lightgray,dimen=inner,fillstyle=hlines*,fillcolor=lightgray,hatchangle=45](0,0)(10,2) 
% Blocos: 
\bloco[linewidth=2pt,linecolor=black]{\xbi}{\ybi}{\wbi}{\hbi}
\bloco[linewidth=2pt,linecolor=black]{\xbii}{\ybii}{\wbii}{\hbii}
% Força: 
\psline[linewidth=3pt,linecolor=blue]{->}(\xfs,\ybi)(\xfe,\ybi) 
% Aceleração: 
\psline[linewidth=3pt,linecolor=blue]{->}(\xas,\ybii)(\xae,\ybii) 
% Texto: 
\rput{0}(\xmi,\ymi){\scalebox{1.5}{$m_1$}}
\rput{0}(\xmii,\ymii){\scalebox{1.5}{$m_2$}}
\rput{0}(\xf,\yf){\scalebox{1.5}{$\mathbf{F}$}}
\rput{0}(\xa,\ya){\scalebox{1.5}{$\mathbf{a}$}}
\end{pspicture}

\end{document}
